%% -*- coding: utf-8 -*-


\section{时态}

\subsection{一般现在时的结构}

\begin{center}
  \begin{talltblr}[
    label = none, % Removes the "Table X:" text
    ]{
    vlines, hlines,
    % 加入支持列表环境的测量参数 (记得在导言区加载 \UseTblrLibrary{varwidth})
    % measure=vbox, stretch=-1,
    rows={halign=c, valign=m},
    columns={halign=c, valign=m},
    column{1}={0.12\linewidth, halign=c, bg=azure9},
    column{2}={0.12\linewidth, halign=c},
    column{3}={0.25\linewidth, halign=l},
    column{4}={0.35\linewidth, halign=l},
    row{1}={font=\large\bfseries, bg=azure9, halign=c}
    }
    类型&句式&结构&举例\\
    \SetCell[r=3]{c}谓语动词是be&肯定句&主语+be（am/is/are）+其他&The new policy is of great significance for environmental protection.\\
        &否定句&主语+be（am/is/are）+not+其他&The project is not (isn't) feasible due to the lack of sufficient funds.\\
        &一般疑问句&be（am/is/are）+主语+其他?&Is the conference room available this afternoon?\\
    \SetCell[r=3]{c}谓语动词是实义动词&肯定句&主语+动词原形/三单+其他&She often participates in international academic seminars.\\
        &否定句&主语+do/does not+动词原形+其他&They don't understand the complexity of the situation.\\
        &一般疑问句&Do/Does+主语+动词原形+其他?&Does he always submit his work on time?\\
  \end{talltblr}
\end{center}


\subsection{一般现在时的标志词：always/frequently/regularly/seldom/sometimes/usually\(\ldots\)}

一般现在时用于表示经常发生的、习惯性的、现在反复出现的动作或状态，常用的时间状语有：always, frequently, regularly, seldom, sometimes, usually, every day, from time to time, now and then, once a week, every other day, as a rule（通常）等。

\subsection{一般现在时的用法}

\begin{center}
  \begin{talltblr}[
    label = none, % Removes the "Table X:" text
    ]{
    vlines, hlines,
    % 加入支持列表环境的测量参数 (记得在导言区加载 \UseTblrLibrary{varwidth})
    % measure=vbox, stretch=-1,
    rows={halign=c, valign=m},
    columns={halign=c, valign=m},
    % column{1}={0.15\linewidth, halign=c, bg=azure9},
    column{1}={0.35\linewidth, halign=l},
    column{2}={0.45\linewidth, halign=l},
    % column{4}={0.35\linewidth, halign=l},
    row{1}={font=\large\bfseries, bg=azure9, halign=c}
    }
    用法&举例\\
    描述客观事实或普遍真理&Light travels faster than sound.\\
    表示人或物的基本特征、特性及目前的状态&The newly bought sofa feels extremely comfortable.\\
    表示当前所发生的动作或存在的状态（这种状态带有一定的持续性）&My parents own a cozy coffee shop in the downtown.\\
    书报的标题、故事的叙述、小说/戏剧/电影等的情节介绍、图片的说明等&Figure 2 illustrates the complete architectural picture.\\
    表示经常发生的、习惯性的、现在反复出现的动作或状态&She usually takes a yoga class on weekends to relieve stress.\\
  \end{talltblr}
\end{center}

\subsection{一般现在时表示将来}

\begin{center}
  \begin{talltblr}[
    label = none, % Removes the "Table X:" text
    ]{
    vlines, hlines,
    % 加入支持列表环境的测量参数 (记得在导言区加载 \UseTblrLibrary{varwidth})
    % measure=vbox, stretch=-1,
    rows={halign=c, valign=m},
    columns={halign=c, valign=m},
    % column{1}={0.15\linewidth, halign=c, bg=azure9},
    column{1}={0.35\linewidth, halign=l},
    column{2}={0.45\linewidth, halign=l},
    % column{4}={0.35\linewidth, halign=l},
    row{1}={font=\large\bfseries, bg=azure9, halign=c}
    }
    用法&举例\\
    在时间状语从句和条件状语从句中，一般现在时表示将来要发生的动作&We'll start the project as soon as all the necessary materials arrive.\\
    谓语动词是come, go, arrive, leave, start, begin return, live等时，一般现在时可以表示将来发生的动作&She goes to Paris for a business trip next week.\\
    一般现在时可以和一个时间短语连用以表示已确定的、对将来的安排&The conference begins on the first day of next month.\\
  \end{talltblr}
\end{center}

\subsection{一般过去时的结构}

\begin{center}
  \begin{talltblr}[
    label = none, % Removes the "Table X:" text
    ]{
    vlines, hlines,
    % 加入支持列表环境的测量参数 (记得在导言区加载 \UseTblrLibrary{varwidth})
    % measure=vbox, stretch=-1,
    rows={halign=c, valign=m},
    columns={halign=c, valign=m},
    column{1}={0.12\linewidth, halign=c, bg=azure9},
    column{2}={0.12\linewidth, halign=l},
    column{3}={0.26\linewidth, halign=l},
    column{4}={0.35\linewidth, halign=l},
    row{1}={font=\large\bfseries, bg=azure9, halign=c}
    }
    类型&句式&结构&举例\\
    \SetCell[r=3]{c}谓语动词是be&肯定句&主语+be（was/were）+其他&She was a teacher before becoming a writer.\\
        &否定句&主语+be（was/were）+not+其他&They weren't satisfied with the results.\\
        &一般疑问句&be（Was/Were）+主语+其他&Were the experimental conditions strictly controlled?\\
    \SetCell[r=3]{c}谓语动词是实义动词&肯定句&主语+动词过去式+其他&The government implemented new policies last year.\\
        &否定句&主语+did not (didn't)+动词原形+其他&The company didn't allocate sufficient funds for product innovation.\\
        &一般疑问句&Did+主语+动词原形+其他&Did the manager approve the proposed marketing strategy?\\
  \end{talltblr}
\end{center}


\subsection{一般过去时的标志词：yesterday/just now/ago/at that time \(\ldots\)}

一般过去时表示过去某个特定时间发生的动作或存在的状态，在一般过去时的句子中，通常都要有表示过去的时间状语，如： yesterday, two days ago, last year, just now, in the old days, the other day, in 1990, at that time, once upon a time, etc.

\subsection{一般过去时的用法}

\begin{center}
  \begin{talltblr}[
    label = none, % Removes the "Table X:" text
    ]{
    vlines, hlines,
    % 加入支持列表环境的测量参数 (记得在导言区加载 \UseTblrLibrary{varwidth})
    % measure=vbox, stretch=-1,
    rows={halign=c, valign=m},
    columns={halign=c, valign=m},
    % column{1}={0.15\linewidth, halign=c, bg=azure9},
    column{1}={0.35\linewidth, halign=l},
    column{2}={0.45\linewidth, halign=l},
    % column{4}={0.35\linewidth, halign=l},
    row{1}={font=\large\bfseries, bg=azure9, halign=c}
    }
    用法&举例\\
    表示过去某个特定时间发生的动作或存在的状态&Did you attend the academic seminar last Friday?\\
    表示过去一段时间内经常发生或反复的动作&She practiced speaking English every day during her stay in the UK.\\
    讲述过去的故事或经历&Once upon a time, a beautiful girl lived in a big castle.\\
    讲述过去连续发生的动作&I woke up early, ate a breakfast, and headed to the exam center.\\
  \end{talltblr}
\end{center}


\subsection{没有时间状语的一般过去时：used to/would\(\ldots\)}

\begin{enumerate}
  \item 描述过去连续发生的动作要用过去时，这种情况下，往往没有表示过去的时间状语，而要通过上下文来理解。比如:The journalist rushed to the scene, interviewed the witnesses and wrote te report.这位记者赶到现场，采访了目击者并撰写了报道。
  \item 用used to do或 would do表示过去经常或反复发生的动作，比如:He used to take awalk in the morning.他以前常常在早晨散步。
  \item 有些情况发生的时间未清楚表明，但实际上是“刚才，刚刚”发生的，属于过去时间，应使用过去时态。常见的表述有“I didn't know/realize.”或“I forgot..”等。比如:I didn't realize the importance of this data until you mentioned it.直到你提到，我才意识到这些数据的重要性。
\end{enumerate}

\subsection{一般将来时will}

\begin{enumerate}
  \item 一般将来时will的结构
        \begin{center}
          \begin{talltblr}[
            label = none, % Removes the "Table X:" text
            ]{
            vlines, hlines,
            % 加入支持列表环境的测量参数 (记得在导言区加载 \UseTblrLibrary{varwidth})
            % measure=vbox, stretch=-1,
            rows={halign=c, valign=m},
            columns={halign=c, valign=m},
            column{1}={0.12\linewidth, halign=c, bg=azure9},
            column{2}={0.3\linewidth, halign=l},
            column{3}={0.43\linewidth, halign=l},
            % column{4}={0.35\linewidth, halign=l},
            row{1}={font=\large\bfseries, bg=azure9, halign=c}
            }
            句式&结构&举例\\
            肯定句&主语+will/shall+动词原形+其他&The government will implement more effective policies.\\
            否定句&主语+will/shall+not（won't）+动词原形+其他&The research team won't complete the experiment within the scheduled time.\\
            一般疑问句&Will/Shall+主语+动词原形+其他&Will the new teaching method improve students' academic performance?\\
          \end{talltblr}
        \end{center}
  \item 一般将来时will的用法
        \begin{center}
          \begin{talltblr}[
            label = none, % Removes the "Table X:" text
            ]{
            vlines, hlines,
            % 加入支持列表环境的测量参数 (记得在导言区加载 \UseTblrLibrary{varwidth})
            % measure=vbox, stretch=-1,
            rows={halign=c, valign=m},
            columns={halign=c, valign=m},
            % column{1}={0.15\linewidth, halign=c, bg=azure9},
            column{1}={0.45\linewidth, halign=l},
            column{2}={0.45\linewidth, halign=l},
            % column{4}={0.35\linewidth, halign=l},
            row{1}={font=\large\bfseries, bg=azure9, halign=c}
            }
            用法&举例\\
            基于个人观点或过去的经验做出预测&I think it will be extremely cold there.\\
            表达目前尚未安排、将来有可能发生的事情&She will probably travel to Europe next year.\\
            表达将来可能发生的非主观的事情或事实&The sun will rise in the east as usual tomorrow.\\
            表达说话者在说话间决定要做的事情，此时常用来提供帮助、作出承诺或给予建议&I will help you with your English if you have any problems.\\
          \end{talltblr}
        \end{center}
\end{enumerate}

\subsection{一般将来时 be going to}

\begin{enumerate}
  \item 一般将来时be going to的结构
        \begin{center}
          \begin{talltblr}[
            label = none, % Removes the "Table X:" text
            ]{
            vlines, hlines,
            % 加入支持列表环境的测量参数 (记得在导言区加载 \UseTblrLibrary{varwidth})
            % measure=vbox, stretch=-1,
            rows={halign=c, valign=m},
            columns={halign=c, valign=m},
            column{1}={0.12\linewidth, halign=c, bg=azure9},
            column{2}={0.3\linewidth, halign=l},
            column{3}={0.43\linewidth, halign=l},
            % column{4}={0.35\linewidth, halign=l},
            row{1}={font=\large\bfseries, bg=azure9, halign=c}
            }
            句式&结构&举例\\
            肯定句&主语+be (am/is/are)+going to+动词原形+其他&He is going to attend an international conference.\\
            否定句&主语+be (am/is/are)+not going to+动词原形+其他&She isn't going to waste her time on meaningless things.\\
            一般疑问句&be (Am/Is/Are)+主语+going to+动词原形+其他&Are you going to study abroad after graduation?\\
          \end{talltblr}
        \end{center}
  \item 一般将来时be going to的用法
        \begin{center}
          \begin{talltblr}[
            label = none, % Removes the "Table X:" text
            ]{
            vlines, hlines,
            % 加入支持列表环境的测量参数 (记得在导言区加载 \UseTblrLibrary{varwidth})
            % measure=vbox, stretch=-1,
            rows={halign=c, valign=m},
            columns={halign=c, valign=m},
            % column{1}={0.15\linewidth, halign=c, bg=azure9},
            column{1}={0.45\linewidth, halign=l},
            column{2}={0.45\linewidth, halign=l},
            % column{4}={0.35\linewidth, halign=l},
            row{1}={font=\large\bfseries, bg=azure9, halign=c}
            }
            用法&举例\\
            表达已经考虑过并打算实施的将来的某种行为&She is going to start a new business project this year.\\
            根据目前事实做出的预测&Look at those dark clouds. It's going to rain.\\
          \end{talltblr}
        \end{center}
\end{enumerate}

\subsection{will和be going to的用法}

\begin{enumerate}
  \item will 和 be going to 跟在 think、doubt、expect、 believe、probably、 certainly、definitely、be sure 这些单词和短语后面，表达了关于将来的某种看法。比如:
        \begin{itemize}
          \item We expect the prices will go down soon.我们期望价格很快会下降。
          \item They are definitely going to buy a new house.他们肯定会买一套新房子。
        \end{itemize}
  \item will和be going to在表达对未来的预测时，虽然都可以表示将来会发生的事情，但在使用上有一些细微差别。will 侧重于基于个人观点或过去的经验做出预测;be going to则侧重于根据目前事实做出预测。比如:
        \begin{itemize}
          \item I'm sure the new movie will be a big hit.我敢肯定这部新电影会大受欢迎。
          \item The water level in the river is rising rapidly, it is going to flood the nearby area.河水水位在迅速上升，附近地区要被淹了。
        \end{itemize}
\end{enumerate}

\subsection{表示预测将来的短语：be likely/bound/certain/estimated/predicted/set to}

除了使用 will 和 be going to，还可以用一些短语来表示将来。此类短语包括 be bound to、be certain to、be estimated to、be likely to、be predicted to、be set to等。比如:
\begin{itemize}
  \item Technological innovations are bound to reshape the future of industries.技术创新必将重塑各行业的未来。
  \item The new policy is set to bring about substantial changes in the housing market.新政策即将给房地产市场带来重大变化。
\end{itemize}

\subsection{一般将来时的标志语：tomorrow/soon/in one week/in the future/eventually\(\ldots\)}

在一般将来时的句子中，通常都会有表示将来的时间状语，如:the day after tomorrow (后天)、tonight (今天晚上)、tomorrow (明天)、this afternoon (今天下午)、next week (下周)、before long (不久后)、in ten years (十年后)、in the future (在将来)、in a little while (过一会儿)、eventually (最终)、sooner or later (迟早)等等。

\subsection{be about to表示将来时}

be about to表示即将发生的动作，结构如下：
\begin{center}
  \begin{talltblr}[
    label = none, % Removes the "Table X:" text
    ]{
    vlines, hlines,
    % 加入支持列表环境的测量参数 (记得在导言区加载 \UseTblrLibrary{varwidth})
    % measure=vbox, stretch=-1,
    rows={halign=c, valign=m},
    columns={halign=c, valign=m},
    column{1}={0.12\linewidth, halign=c, bg=azure9},
    column{2}={0.3\linewidth, halign=l},
    column{3}={0.43\linewidth, halign=l},
    % column{4}={0.35\linewidth, halign=l},
    row{1}={font=\large\bfseries, bg=azure9, halign=c}
    }
    句式&结构&举例\\
    肯定句&主语+be (am/is/are)+about to+动词原形+其他&The plane is about to leave.\\
    否定句&主语+be (am/is/are)+not about to+动词原形+其他&He isn't about to give up hie hard-earned opportunity easily.\\
    一般疑问句&be (Am/Is/Are)+主语+about to+动词原形+其他&Am I about to miss the train?\\
  \end{talltblr}
\end{center}

\subsection{现在进行时的结构}

\begin{center}
  \begin{talltblr}[
    label = none, % Removes the "Table X:" text
    ]{
    vlines, hlines,
    % 加入支持列表环境的测量参数 (记得在导言区加载 \UseTblrLibrary{varwidth})
    % measure=vbox, stretch=-1,
    rows={halign=c, valign=m},
    columns={halign=c, valign=m},
    column{1}={0.12\linewidth, halign=c, bg=azure9},
    column{2}={0.3\linewidth, halign=l},
    column{3}={0.43\linewidth, halign=l},
    % column{4}={0.35\linewidth, halign=l},
    row{1}={font=\large\bfseries, bg=azure9, halign=c}
    }
    句式&结构&举例\\
    肯定句&主语+be (am/is/are)+现在分词+其他&I am analyzing a complex data set.\\
    否定句&主语+be (am/is/are)+not+现在分词+其他&She isn't practicing the piano right now.\\
    一般疑问句&be (Am/Is/Are)+主语+现在分词+其他&Is she preparing for the upcoming exam?\\
  \end{talltblr}
\end{center}

\subsection{现在进行时的标志词：now/at the moment/listen/look\(\ldots\)}

在现在进行时中，经常用到的时间状语有now、at the moment、listen、look 以及一些暗示句子中的动作正在发生的词或句子。比如:
\begin{enumerate}
  \item 当句中出现了now、at the moment等时，表示句子要说明的是现在正在发生的事，动词应用现在进行时。
  \item 当句中出现的时间状语是 these days、this week、this month、this term 等时，如果句子所要表达的意义是在这一阶段正在发生的事，则动词应用现在进行时。
  \item 当句中出现了look、listen等暗示词时，说明后面谓语动词的动作正在发生，动词应用现在进行时。
\end{enumerate}

\subsection{现在进行时的一般用法}

\begin{center}
  \begin{talltblr}[
    label = none, % Removes the "Table X:" text
    ]{
    vlines, hlines,
    % 加入支持列表环境的测量参数 (记得在导言区加载 \UseTblrLibrary{varwidth})
    % measure=vbox, stretch=-1,
    rows={halign=c, valign=m},
    columns={halign=c, valign=m},
    % column{1}={0.15\linewidth, halign=c, bg=azure9},
    column{1}={0.45\linewidth, halign=l},
    column{2}={0.45\linewidth, halign=l},
    % column{4}={0.35\linewidth, halign=l},
    row{1}={font=\large\bfseries, bg=azure9, halign=c}
    }
    用法&举例\\
    表示说话时正在发生的动作&Stop making noise, I'm concentrating on my work.\\
    表示现阶段正在进行的动作，虽然此时此刻该动作不一定正在进行&She is writing a research paper these weeks.\\
    表示此时此刻某一个动作不断地重复&The boy is jumping with great joy at the sight of the toy car.\\
  \end{talltblr}
\end{center}

\subsection{现在进行时的特殊用法}

\begin{center}
  \begin{talltblr}[
    label = none, % Removes the "Table X:" text
    ]{
    vlines, hlines,
    % 加入支持列表环境的测量参数 (记得在导言区加载 \UseTblrLibrary{varwidth})
    % measure=vbox, stretch=-1,
    rows={halign=c, valign=m},
    columns={halign=c, valign=m},
    % column{1}={0.15\linewidth, halign=c, bg=azure9},
    column{1}={0.45\linewidth, halign=l},
    column{2}={0.45\linewidth, halign=l},
    % column{4}={0.35\linewidth, halign=l},
    row{1}={font=\large\bfseries, bg=azure9, halign=c}
    }
    用法&举例\\
    表示动作逐渐变化的过程，常用动词有get, become, come, go等&The problem is becoming more and more complicated.\\
    表示生气、赞叹或厌恶等感情色彩，常与always, constantly, continually, forever等副词连用&He is constantly making the same mistakes.\\
  \end{talltblr}
\end{center}

\subsection{不可用于现在进行时的动词：感官系动词、表示情感、心理状态的动词、拥有等}

并不是所有动词都可以用于“现在进行时”，比如感官系动词、表达情感态度或心理态度的词等通常并不出现在现在进行时的句子里，汇总如下：
\begin{center}
  \begin{talltblr}[
    label = none, % Removes the "Table X:" text
    ]{
    vlines, hlines,
    % 加入支持列表环境的测量参数 (记得在导言区加载 \UseTblrLibrary{varwidth})
    % measure=vbox, stretch=-1,
    rows={halign=c, valign=m},
    columns={halign=c, valign=m},
    column{1}={0.25\linewidth, halign=c, bg=azure9},
    column{2}={0.55\linewidth, halign=l},
    % column{3}={0.35\linewidth, halign=l},
    % column{4}={0.35\linewidth, halign=l},
    row{1}={font=\large\bfseries, bg=azure9, halign=c}
    }
    特殊多次&举例\\
    感官系动词&feel, hear, look, see, small, sound, taste\\
    表示情感态度的动词&adore, despise, dislike, enjoy, feel, hate, like, love, mind, prefer, want\\
    表示心理态度的动词&agree, assume, believe, disagree, forget, hope, know, regret, remember, suppose, think, understand\\
    表示“拥有”的动词&belong, have, own\\
    其他静态动词&appear, seem, contain, resemble (类似、相似), equal\\
  \end{talltblr}
\end{center}

\subsection{现在进行时表示将来}

现在进行时还可以表示按计划或安排即将发生的动作。
\begin{enumerate}
  \item 描述已经确定的计划安排，比如:
        \begin{itemize}
          \item She's having a meeting with her clients this afternoon.她今天下午要和客户开会。
          \item We're starting a new project next month.我们下个月要启动一个新项目。
        \end{itemize}
  \item 描述趋势预测以及流程步骤，比如:
        \begin{itemize}
          \item The raw materials are being transported to the factory next Tuesday.这些原材料将于下周二被运往工厂。
        \end{itemize}
\end{enumerate}

\subsection{将来进行时的结构}

\begin{center}
  \begin{talltblr}[
    label = none, % Removes the "Table X:" text
    ]{
    vlines, hlines,
    % 加入支持列表环境的测量参数 (记得在导言区加载 \UseTblrLibrary{varwidth})
    % measure=vbox, stretch=-1,
    rows={halign=c, valign=m},
    columns={halign=c, valign=m},
    column{1}={0.12\linewidth, halign=c, bg=azure9},
    column{2}={0.3\linewidth, halign=l},
    column{3}={0.43\linewidth, halign=l},
    % column{4}={0.35\linewidth, halign=l},
    row{1}={font=\large\bfseries, bg=azure9, halign=c}
    }
    句式&结构&举例\\
    肯定句&主语+will be+现在分词+其他&She will be attending a conference next Monday.\\
    否定句&主语+will+not+be+现在分词+其他&He won't be working on this project this weekend.\\
    一般疑问句&Will+主语+be+现在分词+其他？&Will they be travelling abroad next month?\\
  \end{talltblr}
\end{center}

\subsection{将来进行时的用法}

\begin{center}
  \begin{talltblr}[
    label = none, % Removes the "Table X:" text
    ]{
    vlines, hlines,
    % 加入支持列表环境的测量参数 (记得在导言区加载 \UseTblrLibrary{varwidth})
    % measure=vbox, stretch=-1,
    rows={halign=c, valign=m},
    columns={halign=c, valign=m},
    % column{1}={0.15\linewidth, halign=c, bg=azure9},
    column{1}={0.45\linewidth, halign=l},
    column{2}={0.45\linewidth, halign=l},
    % column{4}={0.35\linewidth, halign=l},
    row{1}={font=\large\bfseries, bg=azure9, halign=c}
    }
    用法&举例\\
    表示未来某个时刻正在进行或持续的动作&What will you be doing this time next year?\\
    表示按计划或安排将要发生的动作&They will be having a seminar next week.\\
    可用于询问别人的计划、打算（比一般将来时更委婉）&Will you be joining us for dinner?\\
    有时可以代替一般将来时使用，表示一种已经决定获肯定的动作或情况，或表示某动作将继续但还未完成&She will be giving a presentation at the meeting.\\
    与一般将来时连用时，表示将来进行时的动作晚于一般将来时所代表的动作&The construction will end next month and we will be moving into the new building.\\
  \end{talltblr}
\end{center}

\subsection{过去进行时的结构和一般用法}

过去进行时的结构是“was/were+doing”，其用法如下：
\begin{enumerate}
  \item 表示过去某段时间内持续进行的动作或者事情，常用的时间状语有 this morning、all day yesterday、last evening、at that time、when、while 等。比如:
        \begin{itemize}
          \item All day yesterday, the workers were busy renovating the old house.昨天一整天，工人们都在忙着翻新旧房子。
        \end{itemize}
  \item 表示在过去某个时间点发生的事情，比如:What were you doing at 9 p.m. last night?你昨天晚上九点在干什么?
  \item 表示过去某一动作发生时，另一动作正在进行，比如:When the phone rang, I was taking a shower.电话响的时候，我正在洗澡。
\end{enumerate}

\subsection{过去进行时的特殊用法}

过去进行时除了表示“过去正在发生”的事情，还有以下几类用法。
\begin{enumerate}
  \item 与频率副词 always、often、usually等连用，表示过去反复性或习惯性的动作，比如： You were always changing your mind.你总是改变自己的想法。
  \item 表示在过去即将发生的动作，比如:The professor asked the students if they were submitting their dissertations on time.教授问学生们是否会按时提交论文。
  \item 礼貌用法，比如:I was wondering if you could spare a few minutes to review my essay.不知道您能否抽出几分钟帮我看一下我的文章。
  \item 此外，还有一些动词不能用于过去进行时，比如agree、believe、love、meanremember、understand 等。
\end{enumerate}

\subsection{现在完成时的结构}

\begin{center}
  \begin{talltblr}[
    label = none, % Removes the "Table X:" text
    ]{
    vlines, hlines,
    % 加入支持列表环境的测量参数 (记得在导言区加载 \UseTblrLibrary{varwidth})
    % measure=vbox, stretch=-1,
    rows={halign=c, valign=m},
    columns={halign=c, valign=m},
    column{1}={0.12\linewidth, halign=c, bg=azure9},
    column{2}={0.3\linewidth, halign=l},
    column{3}={0.43\linewidth, halign=l},
    % column{4}={0.35\linewidth, halign=l},
    row{1}={font=\large\bfseries, bg=azure9, halign=c}
    }
    句式&结构&举例\\
    肯定句&主语+have/has+过去分词+其他&The researchers have made significant progress.\\
    否定句&主语+have/has+not+过去分词+其他&They haven't decided on the topic of their academic paper.\\
    一般疑问句&Have/Has+主语+过去分词+其他？&Has she participated in any international exchange program before?\\
  \end{talltblr}
\end{center}

\subsection{现在完成时的标志词：since/already/just/yet/never\(\ldots\)}

使用现在完成时的句子中常出现 already、just、yet、never、before、lately、up to now、so far、since 等标志词，其中 since 的用法较为典型，现汇总如下：
\begin{enumerate}
  \item “since +某一确定的时间点 (如具体的时间、日期)”，谓语动词常用现在完成时，且该动词应为延续性动词。比如:She has been committed to her research on AI since last January.从去年一月起，她就一直专注于人工智能方面的研究。
  \item “since+一段时间+ago”，谓语动词常用现在完成时，且该动词应为延续性动词。比如:They have kept in close touch with international scholars since three years ago.从年前开始，他们就一直与国际学者保持密切联系。
  \item “since+从句”，主句谓语动词常用现在完成时，从句谓语动词用一般过去时。比如:The city has experienced rapid development since the new mayor took office. 自从新市长上任以来，这座城市经历了快速发展。
  \item “It is/has been +一段时间+ since从句”,意为“自从……以来已经 (多久)了”。比如:It has been five months since they launched the new project.他们启动这个新项目已经五个月了。
\end{enumerate}

\subsection{现在完成时的用法}

\begin{center}
  \begin{talltblr}[
    label = none, % Removes the "Table X:" text
    ]{
    vlines, hlines,
    % 加入支持列表环境的测量参数 (记得在导言区加载 \UseTblrLibrary{varwidth})
    % measure=vbox, stretch=-1,
    rows={halign=c, valign=m},
    columns={halign=c, valign=m},
    % column{1}={0.15\linewidth, halign=c, bg=azure9},
    column{1}={0.45\linewidth, halign=l},
    column{2}={0.45\linewidth, halign=l},
    % column{4}={0.35\linewidth, halign=l},
    row{1}={font=\large\bfseries, bg=azure9, halign=c}
    }
    用法&举例\\
    表示发生在过去而对现在仍有影响的动作&I have just finished my lunch.\\
    表示从过去某事开始而延续至今的动作或状态（常与for或since等词连用）&I have lived here since 2015.\\
    在句型“This is+adj.最高级+名词+that从句”中，that从句中谓语要用现在完成时&This is the most challenging task that I have ever encountered in my research.\\
    在句型“It is first/second（或其他序数词） time that从句”中，如果强调该动作对现在的影响，that从句中谓语也要用现在完成时&It is the first time that I have been to Shanghai.\\
  \end{talltblr}
\end{center}

\subsection{时间状语在现在完成时中的位置}

\begin{center}
  \begin{talltblr}[
    label = none, % Removes the "Table X:" text
    ]{
    vlines, hlines,
    % 加入支持列表环境的测量参数 (记得在导言区加载 \UseTblrLibrary{varwidth})
    % measure=vbox, stretch=-1,
    rows={halign=c, valign=m},
    columns={halign=c, valign=m},
    column{1}={0.15\linewidth, halign=c, bg=azure9},
    column{2}={0.3\linewidth, halign=l},
    column{3}={0.4\linewidth, halign=l},
    % column{4}={0.35\linewidth, halign=l},
    row{1}={font=\large\bfseries, bg=azure9, halign=c}
    }
    时间状语位置&副词&举例\\
    在助动词和动词之间&already, just, ever, never, recently&She has just received an offer from a prestigious university.\\
    在动词之后&all my life, every day, yet, before, for ages, since 2000, since I was a child, etc.&We haven't finished the project yet.\\
  \end{talltblr}
\end{center}

\subsection{现在完成时vs一般过去时}

\begin{center}
  \begin{talltblr}[
    label = none, % Removes the "Table X:" text
    ]{
    vlines, hlines,
    % 加入支持列表环境的测量参数 (记得在导言区加载 \UseTblrLibrary{varwidth})
    % measure=vbox, stretch=-1,
    rows={halign=c, valign=m},
    columns={halign=c, valign=m},
    column{1}={0.15\linewidth, halign=c, bg=azure9},
    column{2}={0.3\linewidth, halign=l},
    column{3}={0.4\linewidth, halign=l},
    % column{4}={0.35\linewidth, halign=l},
    row{1}={font=\large\bfseries, bg=azure9, halign=c}
    }
    时态&结构&区别\\
    现在完成时&主语+have/has done+其他&表示动作已完成，强调该动作对现在造成的影响或结果\\
    一般过去时&{主语+was/were+其他\\主语+动词过去式+其他}&表示过去某个特定时间发生的动作或存在的状态，与现在没有直接联系\\
  \end{talltblr}
\end{center}

\subsection{过去完成时的结构}

\begin{center}
  \begin{talltblr}[
    label = none, % Removes the "Table X:" text
    ]{
    vlines, hlines,
    % 加入支持列表环境的测量参数 (记得在导言区加载 \UseTblrLibrary{varwidth})
    % measure=vbox, stretch=-1,
    rows={halign=c, valign=m},
    columns={halign=c, valign=m},
    column{1}={0.12\linewidth, halign=c, bg=azure9},
    column{2}={0.3\linewidth, halign=l},
    column{3}={0.43\linewidth, halign=l},
    % column{4}={0.35\linewidth, halign=l},
    row{1}={font=\large\bfseries, bg=azure9, halign=c}
    }
    句式&结构&举例\\
    肯定句&主语+had+过去分词+其他&He had completed all the assigned tasks before the deadline arrived.\\
    否定句&主语+had+not+过去分词+其他&They had not made any decision regarding the investment plan until last week.\\
    一般疑问句&Had+主语+过去分词+其他？&Had the researchers obtained sufficient data before they began the analysis?\\
  \end{talltblr}
\end{center}

\subsection{过去完成时的用法}

\begin{center}
  \begin{talltblr}[
    label = none, % Removes the "Table X:" text
    ]{
    vlines, hlines,
    % 加入支持列表环境的测量参数 (记得在导言区加载 \UseTblrLibrary{varwidth})
    % measure=vbox, stretch=-1,
    rows={halign=c, valign=m},
    columns={halign=c, valign=m},
    % column{1}={0.15\linewidth, halign=c, bg=azure9},
    column{1}={0.45\linewidth, halign=l},
    column{2}={0.45\linewidth, halign=l},
    % column{4}={0.35\linewidth, halign=l},
    row{1}={font=\large\bfseries, bg=azure9, halign=c}
    }
    用法&举例\\
    表示较早的过去，即某一时刻之前已完成的动作或状态&By the end of last month, the research team had collected all the necessary data.到上个月月底,研究团队已经收集到了所有必要的数据。\\
    表示由过去的某一时刻开始，一直延续到过去另一时间点的动作或状态 (常与for、since 等连用)&He said he had worked in that company since 2015.他说自2015年以来他就在那家公司工作。\\
    在含有定语从句的复合句中，如果叙述的是过去发生的事，先发生的动作要用过去完成时&The professor praised the students who had made remarkable progress.教授表扬了那些取得显著进步的学生。\\
    用于宾语从句:如果主句中用了know、realize、think、suppose、find、decide 等动词的一般过去式，且宾语从句中的动作先于主句的动作，通常用过去完成时&They realized they had underestimated the difficulty of the task.他们意识到他们低估了任务的难度。\\
    用来表示未曾实现的愿望或打算，此时谓语动词经常是hope、want、wish、mean、intend、attempt、think、expe等，表示“原本……、未能……”等语气&She had hoped to get a high score on the exam, but she didn't perform well.\\
  \end{talltblr}
\end{center}

\subsection{将来完成时的结构与用法}

\begin{enumerate}
  \item 将来完成时的结构
        \begin{center}
          \begin{talltblr}[
            label = none, % Removes the "Table X:" text
            ]{
            vlines, hlines,
            % 加入支持列表环境的测量参数 (记得在导言区加载 \UseTblrLibrary{varwidth})
            % measure=vbox, stretch=-1,
            rows={halign=c, valign=m},
            columns={halign=c, valign=m},
            column{1}={0.12\linewidth, halign=c, bg=azure9},
            column{2}={0.3\linewidth, halign=l},
            column{3}={0.43\linewidth, halign=l},
            % column{4}={0.35\linewidth, halign=l},
            row{1}={font=\large\bfseries, bg=azure9, halign=c}
            }
            句式&结构&举例\\
            肯定句&主语+will have+过去分词+其他&By next month, I will have finished my research.\\
            否定句&主语+will not have+过去分词+其他&We will not have finished analyzing the market data by next week.\\
            一般疑问句&Will+主语+have+过去分词+其他？&Will you have submitted your application by the deadline?\\
          \end{talltblr}
        \end{center}
  \item 将来完成时的用法
        \begin{center}
          \begin{talltblr}[
            label = none, % Removes the "Table X:" text
            ]{
            vlines, hlines,
            % 加入支持列表环境的测量参数 (记得在导言区加载 \UseTblrLibrary{varwidth})
            % measure=vbox, stretch=-1,
            rows={halign=c, valign=m},
            columns={halign=c, valign=m},
            % column{1}={0.15\linewidth, halign=c, bg=azure9},
            column{1}={0.25\linewidth, halign=l},
            column{2}={0.35\linewidth, halign=l},
            column{3}={0.3\linewidth, halign=l},
            row{1}={font=\large\bfseries, bg=azure9, halign=c}
            }
            用法&举例&标志词\\
            表示在将来某个时间之前完成的动作，并且往往对将来某一时间产生影响&We will have been married for 5 years by next month.&\SetCell[r=2]{l}将来完成时常用时间状语为“by+将来某个时间”；再使用将来完成时的句子中，by the time引导的时间状语从句中使用一般现在时\\
            表示根据某情况做出的对未来的推测，相当有must have done&They will have reached home by now considering their early start.&\\
          \end{talltblr}
        \end{center}
\end{enumerate}

\subsection{过去将来时的结构与用法}

\begin{enumerate}
  \item 过去将来时的结构
        \begin{center}
          \begin{talltblr}[
            label = none, % Removes the "Table X:" text
            ]{
            vlines, hlines,
            % 加入支持列表环境的测量参数 (记得在导言区加载 \UseTblrLibrary{varwidth})
            % measure=vbox, stretch=-1,
            rows={halign=c, valign=m},
            columns={halign=c, valign=m},
            column{1}={0.12\linewidth, halign=c, bg=azure9},
            column{2}={0.4\linewidth, halign=l},
            column{3}={0.33\linewidth, halign=l},
            % column{4}={0.35\linewidth, halign=l},
            row{1}={font=\large\bfseries, bg=azure9, halign=c}
            }
            句式&结构&举例\\
            肯定句&{主语+would+动词原形+其他\\主语+be (was/were)+going to +动词原形+其他}&She was going to visit her grandparents that weekend.\\
            否定句&{主语+would+not+动词原形+其他\\主语+be (was/were)+not+going to +动词原形+其他}&They said they would not attend the meeting.\\
            一般疑问句&{Would+主语+动词原形+其他\\be (was/were)+主语+going to +动词原形+其他}？&Was she going to study abroad?\\
          \end{talltblr}
        \end{center}
  \item 过去将来时的用法
        \begin{center}
          \begin{talltblr}[
            label = none, % Removes the "Table X:" text
            ]{
            vlines, hlines,
            % 加入支持列表环境的测量参数 (记得在导言区加载 \UseTblrLibrary{varwidth})
            % measure=vbox, stretch=-1,
            rows={halign=c, valign=m},
            columns={halign=c, valign=m},
            % column{1}={0.15\linewidth, halign=c, bg=azure9},
            column{1}={0.35\linewidth, halign=l},
            column{2}={0.5\linewidth, halign=l},
            row{1}={font=\large\bfseries, bg=azure9, halign=c}
            }
            用法&举例\\
            表示从过去某一时间看将要发生的动作或存在的状态&The report predicted that the economic situation would improve in the following year.报告预测经济形势在接下来的-年内会改善。\\
            表示过去的打算、计划或意图&They were going to have a picnic last Sunday.他们上周日原本打算去野餐。\\
          \end{talltblr}
        \end{center}
\end{enumerate}

\subsection{现在完成进行时的结构与用法}

\begin{enumerate}
  \item 现在完成进行时的结构
        \begin{center}
          \begin{talltblr}[
            label = none, % Removes the "Table X:" text
            ]{
            vlines, hlines,
            % 加入支持列表环境的测量参数 (记得在导言区加载 \UseTblrLibrary{varwidth})
            % measure=vbox, stretch=-1,
            rows={halign=c, valign=m},
            columns={halign=c, valign=m},
            column{1}={0.12\linewidth, halign=c, bg=azure9},
            column{2}={0.4\linewidth, halign=l},
            column{3}={0.33\linewidth, halign=l},
            % column{4}={0.35\linewidth, halign=l},
            row{1}={font=\large\bfseries, bg=azure9, halign=c}
            }
            句式&结构&举例\\
            肯定句&主语+have/has+been+现在分词+其他&The researcher has been conducting experiments on this new material for months.\\
            否定句&主语+have/has+not+been+现在分词+其他&He hasn't been teaching here these years.这些年他并没有一直在这儿教书。\\
            一般疑问句&have/has+主语+been+现在分词+其他？&Have you been studying for the math test today?你今天一直在准备数学测试吗?\\
          \end{talltblr}
        \end{center}
  \item 现在完成进行时的用法
        \begin{center}
          \begin{talltblr}[
            label = none, % Removes the "Table X:" text
            ]{
            vlines, hlines,
            % 加入支持列表环境的测量参数 (记得在导言区加载 \UseTblrLibrary{varwidth})
            % measure=vbox, stretch=-1,
            rows={halign=c, valign=m},
            columns={halign=c, valign=m},
            % column{1}={0.15\linewidth, halign=c, bg=azure9},
            column{1}={0.35\linewidth, halign=l},
            column{2}={0.5\linewidth, halign=l},
            row{1}={font=\large\bfseries, bg=azure9, halign=c}
            }
            用法&举例\\
            表示动作从过去某个时刻开始一直持续到现在，并有可能继续下去，强调进行的过程 (常和 for 或 since 引导的时间状语连用)&The environmentalist has been advocating for wildlife protection for a decade.这位环保主义者十年来一直在倡导野生动物保护。\\
            表示动作在一段持续的时间内多次重复&Bob has been phoning Nancy every night for the past week.这一星期以来鲍勃每天晚上都给南希打电话。\\
          \end{talltblr}
        \end{center}
\end{enumerate}

\subsection{现在完成进行时vs现在完成时}

\begin{center}
  \begin{talltblr}[
    label = none, % Removes the "Table X:" text
    ]{
    vlines, hlines,
    % 加入支持列表环境的测量参数 (记得在导言区加载 \UseTblrLibrary{varwidth})
    measure=vbox, stretch=-1,
    rows={halign=c, valign=m},
    columns={halign=c, valign=m},
    % column{1}={0.15\linewidth, halign=c, bg=azure9},
    % column{2}={0.35\linewidth, halign=l},
    column{3}={0.45\linewidth, halign=l},
    % column{4}={0.35\linewidth, halign=l},
    row{1}={font=\large\bfseries, bg=azure9, halign=c}
    }
    时态&结构&区别\\
    现在完成时&have/has done&{
                              \begin{enumerate}[nosep]
                                \item 表示动作已完成
                                \item 强调动作的结果
                              \end{enumerate}
                              }\\
    现在完成进行时&have/has been doing&{
                                        \begin{enumerate}[nosep]
                                          \item 表示动作在迄今为止的一段时间内持续进行或目前仍在进行
                                          \item 强调动作持续进行的状态
                                        \end{enumerate}
                                        }\\
  \end{talltblr}
\end{center}

\subsection{使用完成时态的句型}

\begin{enumerate}
  \item 在句型“It/This/That is+形容词最高级+that从句”中,that从句中谓语要用现在完成时比如:It is the most interesting book that I have ever read.这是我读过的最有趣的书。
  \item 在句型“It/This/That is+序数词+that 从句”中，如果强调该动作对现在的影响，that从句中谓语也要用现在完成时。比如:It is the third time that she has won the award. 这是她第三次获得这个奖项了。
\end{enumerate}

\subsection{过去完成进行时的结构}

\begin{center}
  \begin{talltblr}[
    label = none, % Removes the "Table X:" text
    ]{
    vlines, hlines,
    % 加入支持列表环境的测量参数 (记得在导言区加载 \UseTblrLibrary{varwidth})
    % measure=vbox, stretch=-1,
    rows={halign=c, valign=m},
    columns={halign=c, valign=m},
    column{1}={0.12\linewidth, halign=c, bg=azure9},
    column{2}={0.4\linewidth, halign=l},
    column{3}={0.33\linewidth, halign=l},
    % column{4}={0.35\linewidth, halign=l},
    row{1}={font=\large\bfseries, bg=azure9, halign=c}
    }
    句式&结构&举例\\
    肯定句&主语+had+been+现在分词+其他&\SetCell[r=3]{l}The architect had been designing the new museum for three years before it was approved.\\
    否定句&主语+had+not+been+现在分词+其他&\\
    一般疑问句&Had+主语+been+现在分词+其他？&\\
  \end{talltblr}
\end{center}


\subsection{过去完成进行时的用法}

\begin{center}
  \begin{talltblr}[
    label = none, % Removes the "Table X:" text
    ]{
    vlines, hlines,
    % 加入支持列表环境的测量参数 (记得在导言区加载 \UseTblrLibrary{varwidth})
    % measure=vbox, stretch=-1,
    rows={halign=c, valign=m},
    columns={halign=c, valign=m},
    % column{1}={0.15\linewidth, halign=c, bg=azure9},
    column{1}={0.35\linewidth, halign=l},
    column{2}={0.5\linewidth, halign=l},
    row{1}={font=\large\bfseries, bg=azure9, halign=c}
    }
    用法&举例\\
    表示过去某一时刻之前开始的动作或产生的状态一直延续到过去某一时刻&When the meeting started, I had been preparing the materials for an hour.会议开始时，我已经准备了一个小时的材料。\\
    用在某些从句中，表示从句的动作发在主句的动作之前而且对其有影响&I heard you had been looking for me.我听说你一直在找我。\\
  \end{talltblr}
\end{center}


%%% Local Variables:
%%% TeX-engine: xetex
%%% TeX-master: "../IELTS_300_Sentences"
%%% End:

% LocalWords:  vlines hlines halign valign xetex LocalWords
